% Front matter — versão PT-PT. Espelho de ../thesis/frontmatter/frontmatter.tex (mesmo conteúdo).

\frontmatter % numeração romana (i, ii, iii, ...) nas páginas pré-conteúdo

\pagestyle{plain}

%----------------------------------------------------------------------------------------
%	PÁGINA DE TÍTULO
%----------------------------------------------------------------------------------------

\maketitlepage

%----------------------------------------------------------------------------------------
%	DECLARAÇÃO DE INTEGRIDADE
%----------------------------------------------------------------------------------------

\begin{soi}

\vspace{1cm}

Declaro que realizei este trabalho académico com integridade.

Não plagiei nem apliquei qualquer forma de uso indevido de informação ou de falsificação de resultados
ao longo do processo que conduziu à sua elaboração.

Assim, o trabalho apresentado neste documento é original e da minha autoria, não tendo sido anteriormente
utilizado para qualquer outro fim.

Declaro ainda que tomei pleno conhecimento do Código de Conduta Ética do P.PORTO.

\vspace{0.5cm}

\textbf{Uso de Inteligência Artificial.} Em conformidade com as orientações do P.PORTO/ISEP, declaro que
foram utilizadas ferramentas de inteligência artificial generativa ao longo do desenvolvimento deste
trabalho, e declaro-o na sua extensão real e não de forma minimizada. Foram usadas para escrever e
reestruturar uma parte substancial do código-fonte do sistema e dos seus testes automáticos, para
implementar os procedimentos de avaliação, para redigir e editar prosa neste documento e nos materiais
de apoio, e para conduzir revisões que encontraram defeitos tanto no software como no texto.

O que é meu é o trabalho para onde essas ferramentas foram apontadas. A escolha do problema e do
público, as questões de investigação, as restrições que definem o sistema (apenas fontes de dados
gratuitas, explicabilidade primeiro, e a recusa de prever preços), e todas as decisões sobre o que
construir, promover, manter, estreitar ou descartar são minhas. Onde uma medição contrariou uma
afirmação que eu tinha feito, a afirmação foi retirada ou estreitada em vez de defendida; essas
retractações estão registadas no Apêndice~\ref{app:evidence_matrix}. Todos os valores reportados são
produzidos por um procedimento versionado que pode ser re-corrido. Revi o conteúdo desta dissertação, e
ele, os erros que nele restem, e a sua defesa são da minha responsabilidade.

\vspace{1cm}

ISEP, Porto, \today % NOTA: substituir pela data de entrega; confirmar a redação exata exigida pela MEIA/ISEP.

\end{soi}

%----------------------------------------------------------------------------------------
%	DEDICATÓRIA (opcional)
%----------------------------------------------------------------------------------------
\begin{dedicatory}
% A dedicatória é tua. Está aqui a versão contida; a alternativa mais calorosa está na linha
% comentada abaixo — troca uma pela outra se preferires, ou escreve a tua.
À minha família.
% À minha família, pela paciência que isto também lhes custou.
\end{dedicatory}

%----------------------------------------------------------------------------------------
%	RESUMO (língua principal: Português)
%----------------------------------------------------------------------------------------

\begin{abstract}

% Resumo em Português (obrigatório para trabalho em Inglês). Coerente com o Abstract.

Esta dissertação apresenta o InvestiGator, um sistema de alertas financeiros explicável (XAI-first)
para investidores de retalho no mercado acionista dos Estados Unidos. Envia alertas via Telegram a
partir de dois gatilhos independentes, movimentos abruptos de mercado detetados como anomalias
estatísticas e novas notícias financeiras, cada um com uma explicação rastreável: o evento, o
raciocínio, as fontes e os precedentes históricos. O seu núcleo é um motor de correlação
notícia--mercado que recupera notícias passadas análogas e mede o impacto que esses precedentes
tiveram, à maneira de um estudo de evento e sem lookahead. Uma camada generativa sintetiza estes
resultados em prosa, com cada afirmação ligada à evidência que cita e verificada antes da entrega.
Enquanto contribuição de Engenharia de Inteligência Artificial, o trabalho integra e avalia
criticamente componentes existentes num sistema funcional, explicável e reproduzível. Em dados
reais, a recuperação semântica superou claramente as baselines lexical e triviais (precision@5
cross-ticker de $0{,}51$, face a $0{,}35$ lexical e $0{,}24$ do acaso) e o detetor normalizado pela
volatilidade disparou de forma muito mais consistente entre ações do que um limiar fixo. Um modelo
de triagem de materialidade treinado em $79\,753$ exemplos de vários anos melhorou $1.67\times$ a
precisão dos alertas dentro do orçamento diário em dados retidos, mas nenhum modelo com texto bateu
a baseline de só-volatilidade, e em produção o mesmo gate não mostrou benefício significativo sobre
a população, distribuída de outra forma, que de facto vê. Ambos os resultados são reportados tal
como são, com as limitações explícitas.

\end{abstract}

\begin{abstractotherlanguage}

% Abstract em Inglês (língua principal). <= 200 palavras; alinhado com as conclusões.

This dissertation presents InvestiGator, an explainable (XAI-first) financial-alert system for retail
investors in the US equity market. It raises Telegram alerts from two independent triggers, abrupt
market movements detected as statistical anomalies and incoming financial news, each with
a traceable explanation: the event, the reasoning, the sources and historical precedents. Its core is
a news--market correlation engine that retrieves analogous past news and measures the impact those
precedents had, event-study style and without lookahead. A generative layer synthesises
these results into prose, each claim bound to the evidence it cites and verified before delivery. As
an Artificial Intelligence Engineering contribution, the work integrates and critically evaluates
existing components in a functional, explainable, reproducible system. On real data,
semantic retrieval clearly outperformed lexical and trivial baselines (cross-ticker precision@5 of
$0.51$ against $0.35$ lexical and $0.24$ random), and the volatility-normalised detector fired far
more consistently across stocks than a fixed threshold. A materiality-triage model trained on
$79{,}753$ multi-year examples raised within-budget alert precision $1.67$-fold on held-out data, yet
no text model beat the volatility-only baseline, and in deployment the same gate showed no
significant benefit on the differently distributed population it sees. Both are reported as
they stand, with limitations stated explicitly.

\end{abstractotherlanguage}

%----------------------------------------------------------------------------------------
%	AGRADECIMENTOS (opcional)
%----------------------------------------------------------------------------------------

\begin{acknowledgements}
% RASCUNHO para o aluno rever e reescrever na sua voz. A gratidão é dele; o texto abaixo é um
% ponto de partida com os factos certos (nomes, papéis, e o que cada um contribuiu), não a
% versão final. Verificar: grafia "Sistrade"; se quiser nomear colegas, acrescentar aqui.
Ao meu orientador, Prof.\ Luís Gomes, e ao meu coorientador, Rafael Silva, por todo o apoio ao longo
desta tese. Por trocarem ideias comigo e me ajudarem a decidir a direcção do tema; por me mostrarem
exemplos de aplicações que já existem e me manterem a par do que se vai fazendo na área; e por me
guiarem na elaboração deste documento. Muito do que aqui está tomou a forma que tem porque houve
antes uma conversa com eles.

À minha família, por todo o apoio e por acreditarem em mim. Por me incentivarem sempre, e pela
alegria com que foram acompanhando a construção disto — que é, para mim, a melhor parte de o ter
feito.

Aos meus colegas da Sistrade, e à Sistrade, pela flexibilidade que me permitiu seguir este mestrado
a par do trabalho, pelas opiniões e pelo feedback que foram dando sobre a aplicação, e por serem boa
companhia todos os dias.
% ⚠️ NÃO acrescentar aqui um agradecimento a participantes de estudo: o estudo de utilidade
%    ainda não foi corrido, e agradecer a pessoas que não existem seria fabricação — numa
%    secção onde ninguém iria verificar. Se o estudo for corrido antes da entrega, este é o
%    sítio certo para lhes agradecer, e aí passa a ser verdade.
\end{acknowledgements}

%----------------------------------------------------------------------------------------
%	ÍNDICES
%----------------------------------------------------------------------------------------

\tableofcontents

\listoffigures

\listoftables

\iflanguage{portuguese}{
\renewcommand{\listalgorithmname}{Lista de Algor\'itmos}
}
\listofalgorithms
\addchaptertocentry{\listalgorithmname}

% Lista de excertos de código. Reposta em 2026-08: o repositório não acompanha a dissertação,
% por isso os pontos onde a garantia é *feita* têm de ser legíveis no próprio documento.
\renewcommand{\lstlistlistingname}{Lista de Excertos de C\'odigo}
\lstlistoflistings
\addchaptertocentry{\lstlistlistingname}

%----------------------------------------------------------------------------------------
%	SÍMBOLOS
%----------------------------------------------------------------------------------------

\begin{symbols}{lll} % Lista de símbolos (três colunas)

$r_t$ & retorno logarítmico de um ativo no dia $t$ & --- \\
$\mu_t$ & média móvel dos retornos & --- \\
$\sigma_t$ & desvio-padrão móvel dos retornos & --- \\
$z_t$ & z-score do retorno no dia $t$ & --- \\
$k$ & limiar de anomalia (em desvios-padrão) & --- \\
$w$ & comprimento da janela móvel & \si{\day} \\

\end{symbols}

%----------------------------------------------------------------------------------------
%	ACRÓNIMOS
%----------------------------------------------------------------------------------------

\newcommand{\listacronymname}{List of Acronyms}
\iflanguage{portuguese}{
\renewcommand{\listacronymname}{Lista de Acr\'onimos}
}

\glsresetall

\printnoidxglossary[title=\listacronymname,type=\acronymtype,style=long]

%----------------------------------------------------------------------------------------
%	FIM DO FRONT MATTER
%----------------------------------------------------------------------------------------

\mainmatter
\pagestyle{thesis}
